\input{preamble}

\begin{document}

\title{Art and Culture}
\maketitle

\phantomsection
\label{section-phantom}

\tableofcontents

\section{Abstração e arte brasileira no pós-guerra}

\subsection{Um mapa não linear}

A história da arte brasileira do pós-guerra não deve ser entendida como uma
sequência linear em que figuração, abstração, concretismo e neoconcretismo se
substituem. Em particular, ``figuração'' não é o nome de um primeiro movimento:
é uma categoria ampla para obras que representam figuras, objetos ou espaços
reconhecíveis.

Nos anos 1940--1960 coexistem diferentes vertentes abstratas. Duas famílias
úteis para organizar o período são a abstração informal e a abstração
geométrica. A primeira privilegia gesto, matéria, ritmo, cor e subjetividade; a
segunda privilegia relações formais e construção geométrica. Dentro desta
última se desenvolve o concretismo brasileiro nos anos 1950 e, em diálogo e
ruptura com ele, o neoconcretismo a partir de 1959.

\subsection{Concretismo e neoconcretismo}

A arte concreta busca uma linguagem visual autônoma, fundada na ordem, na
racionalidade e na construção objetiva da forma. No Brasil, um núcleo decisivo
foi o Grupo Ruptura, em São Paulo. No Rio de Janeiro, o Grupo Frente teve Ivan
Serpa como figura central e uma prática mais heterodoxa dentro do campo
geométrico.

O neoconcretismo reage ao caráter excessivamente rígido e racionalista do
programa concreto. Conserva a investigação geométrica, mas introduz uma relação
mais sensível, fenomenológica, espacial e participativa com a obra. Lygia
Clark, Lygia Pape e Hélio Oiticica são nomes centrais desta passagem. Uma
fórmula útil é: no concretismo, geometria construída; no neoconcretismo,
geometria experimentada.

Ivan Serpa é particularmente importante como figura de ligação. Além de sua
produção, sua atividade docente no MAM Rio ajudou a formar o ambiente artístico
carioca do qual emergiram experiências neoconcretas e posteriores.

Artistas como Alfredo Volpi, Maria Leontina e Milton Dacosta dialogam fortemente
com a abstração geométrica do período sem se reduzirem simplesmente aos
programas doutrinários concreto ou neoconcreto. Em Volpi, por exemplo, a
geometria das fachadas e bandeirinhas conserva uma qualidade artesanal e
pictórica muito própria.

\subsection{Abstração informal}

Contemporaneamente ao concretismo, desenvolve-se no Brasil uma abstração
informal. Aqui a tela é entendida como campo de investigação do próprio ato de
pintar: gesto, manchas, tensões, matéria e cor ganham autonomia. Essa família
tem paralelos internacionais no expressionismo abstrato norte-americano e no
Art Informel europeu.

Tomie Ohtake é uma referência importante. Sua produção dos anos 1950--1960
explora gesto, matéria e grandes relações cromáticas, em contraste com a
precisão programática da arte concreta. Algumas obras, como sua \emph{Pintura}
de 1964, podem lembrar visualmente os campos de cor de Mark Rothko, embora a
superfície de Ohtake preserve uma presença própria do gesto e da matéria.

\subsection{Genealogia internacional da abstração geométrica}

O concretismo brasileiro pertence a uma constelação internacional mais ampla de
arte abstrata e construtiva. Uma genealogia simplificada passa pelo cubismo do
início do século XX e, mais diretamente, por Malevich, Mondrian, o
construtivismo russo, De Stijl e a Bauhaus. Nos anos 1950--1960 surgem ou se
consolidam diferentes respostas locais: concretismo e neoconcretismo no Brasil,
Op Art na Europa, hard-edge painting nos Estados Unidos e tradições
construtivas importantes na Argentina e no Uruguai.

Piet Mondrian é um ancestral particularmente direto. Depois do contato com o
cubismo em Paris, sua pintura passa gradualmente da representação de árvores e
paisagens à organização autônoma de linhas, planos e cores. No neoplasticismo,
verticais e horizontais, cores primárias e relações assimétricas de equilíbrio
constituem a própria linguagem da pintura. Sua obra mostra como o cubismo é uma
das grandes rupturas que tornam possível a abstração geométrica, embora esta
não possa ser reduzida simplesmente a uma derivação do cubismo.

\subsection{Brasil e México}

Há afinidades fortes entre a abstração geométrica brasileira e a arte mexicana
do pós-guerra, mas não existe uma correspondência simples entre concretismo e
neoconcretismo brasileiros e movimentos mexicanos de mesmo tipo. No México,
a chamada Generación de la Ruptura e artistas como Vicente Rojo, Mathias
Goeritz e Gunther Gerzso participam da abertura para linguagens abstratas
internacionais, num contexto marcado também pela reação à hegemonia da Escola
Mexicana de Pintura e do muralismo.

Vicente Rojo (1932--2021), ativo no México sobretudo a partir dos anos 1950,
tem afinidades visuais com a abstração geométrica brasileira: construção em
séries, módulos, repetição, cor e uma relação importante com o design gráfico.
Maria Leontina e Milton Dacosta oferecem comparações particularmente naturais.
A relação não deve, porém, ser confundida automaticamente com contato pessoal
entre esses artistas.

Um vínculo histórico México--Brasil mais concreto aparece na poesia concreta.
O grupo brasileiro Noigandres participa de uma circulação internacional que
chega ao México; Mathias Goeritz organiza em 1966, na UNAM, a exposição
\emph{Poesía Concreta Internacional}, na qual Vicente Rojo participa. Assim,
geometria, tipografia, poesia, livro e design constituem uma zona real de
intercâmbio entre os dois países.

\subsection{Retorno da figuração e arte experimental}

Nos anos 1960--1980, a figuração reaparece com novas funções. A Nova Figuração
brasileira incorpora cultura de massa, imprensa, televisão e referências
cotidianas, e no contexto da ditadura militar torna-se também veículo de
crítica social e política. Não se trata de retornar à figuração anterior à
abstração: a imagem retorna depois do cubismo, da abstração e da cultura de
massa.

Ao mesmo tempo, a arte experimental e conceitual radicaliza outra pergunta. Em
vez de buscar apenas uma nova linguagem pictórica, questiona o próprio objeto
artístico, seus suportes, sua circulação e as instituições que o legitimam. A
ênfase pode deslocar-se do objeto para o conceito ou procedimento. Performance,
fotografia, intervenção urbana, corpo e ready-made passam a integrar o campo da
arte. No Brasil, esse desenvolvimento adquire uma dimensão política específica
sob a ditadura.

\subsection{Depois de 1990}

A partir dos anos 1990 torna-se menos convincente narrar a arte brasileira como
uma sucessão de movimentos estilísticos dominantes. Figuração, abstração,
pintura, práticas conceituais, performance, vídeo e meios digitais coexistem.
A revisão dos cânones, a ampliação dos agentes reconhecidos pelo sistema da arte
e questões de memória, território, corpo, identidade e pertencimento tornam-se
centrais. A ideia moderna de ruptura continua presente, mas já não produz uma
sequência tão clara de novos ``ismos''.

\end{document}
